% Eigenständiger Prosabeweis; alle nummerierten Aussagen bleiben in Band 46.
\newcommand{\FranklProofText}[2]{%
  \href{\FranklEditionDirectory_B46-frankl-proofs.pdf\#frankl.proof.#1}{#2}%
}
\section*{Eine kleine Familie und eine auffällige Häufigkeit}
\hypertarget{frankl.reading.example}{}

Wir beginnen mit fünf Mengen. Für drei verschiedene Elemente
\(a,b,c\) betrachten wir die Familie
\[
  M=\{\varnothing,\{a,b\},\{a,c\},\{b,c\},\{a,b,c\}\}.
\]
Eine Familie ist hier selbst eine Menge, deren Elemente Mengen
sind. Wir zählen also fünf verschiedene Mengen; mehrfaches
Aufschreiben derselben Menge würde die Familie nicht vergrößern.
Die Elemente \(a,b,c\) und die fünf Mengen, in denen sie
vorkommen können, übernehmen verschiedene Rollen.

Unsere Familie hat eine besondere Eigenschaft: Vereinigen wir
zwei ihrer Mengen, so gehört das Ergebnis wieder zu ihr.
Mit \(\varnothing\) vereinigt ändert sich nichts. Mit
\(\{a,b,c\}\) vereinigt entsteht stets \(\{a,b,c\}\).
Zwei gleiche Zweiermengen ergeben wieder diese Zweiermenge,
zwei verschiedene Zweiermengen hingegen \(\{a,b,c\}\).
Damit sind alle Möglichkeiten erfasst. Man nennt eine Familie
mit dieser Eigenschaft \emph{vereinigungsabgeschlossen}.

Nun zählen wir, wie oft jedes der drei Elemente vorkommt.

\FranklIncidenceTable

Jedes Element liegt in drei der fünf Mengen, also in mehr als
der Hälfte. Ist das eine Besonderheit unseres symmetrischen
Beispiels? Oder erzwingt der Vereinigungsabschluss allgemein,
dass wenigstens ein Element in mindestens der Hälfte vorkommt?
Diese Frage führt zu Frankls Vermutung.

Der folgende Beweis löst die Frage für zwei besondere Fälle
vollständig: wenn die Familie eine Einermenge oder eine
Zweiermenge als Mitglied enthält. Der zweite Fall umfasst
unser Beispiel, obwohl dessen Träger \(\bigcup M=\{a,b,c\}\)
drei Elemente hat. Entscheidend ist die Größe eines
\emph{Mitglieds}, nicht die Größe des gesamten Trägers.

\par\Needspace{12\baselineskip}
\section*{Was \glqq mindestens die Hälfte\grqq{} bedeutet}
\hypertarget{frankl.reading.frequency}{}

Von nun an sei \(M\) eine endliche Familie von Mengen.
\emph{Endlich} bezieht sich dabei auf die Anzahl ihrer Mitglieder.
Die einzelnen Mengen dürfen in unseren Beweisen auch unendlich
sein. Für ein Element \(x\) zerlegen wir die Familie in
\[
  M_x^+=\{X\in M\mid x\in X\},
  \qquad
  M_x^-=\{X\in M\mid x\notin X\}.
\]
In Band~46 heißen diese beiden Teilfamilien
\(\FamContaining{M}{x}\) und \(\FamAvoiding{M}{x}\).
Jedes Mitglied von \(M\) gehört zu genau einer von ihnen:
Es enthält \(x\) oder enthält \(x\) nicht. Deshalb gilt
\[
  |M|=|M_x^+|+|M_x^-|.
\]
Die senkrechten Striche bezeichnen jeweils die Anzahl der
Mengen in der angegebenen endlichen Familie.

Wir nennen \(x\) \emph{halbhäufig}, wenn es zum Träger
\(\bigcup M\) gehört und in mindestens der Hälfte der
Mitglieder vorkommt. Die entscheidende Umformung lautet
\begin{equation}
  |M_x^+|\geq\frac{|M|}{2}
  \quad\Longleftrightarrow\quad
  |M_x^-|\leq|M_x^+|.
  \tag{H}\label{frankl-reading-half}
\end{equation}
Denn nach der Zerlegung bedeutet die linke Seite
\(2|M_x^+|\geq|M_x^+|+|M_x^-|\); nach Abziehen
von \(|M_x^+|\) bleibt genau die rechte Seite.
Bei fünf Mengen sind somit mindestens drei Vorkommen nötig.

Um diesen Vergleich zu beweisen, genügt eine \emph{injektive
Abbildung} von \(M_x^-\) nach \(M_x^+\). Sie weist jeder
Menge ohne \(x\) eine Menge mit \(x\) zu, wobei verschiedene
Ausgangsmengen verschiedene Bilder bekommen. Die Zielklasse
muss dann mindestens so viele Mitglieder haben wie die
Ausgangsklasse. Nicht jede Menge mit \(x\) muss dabei als
Bild auftreten. Gerade die möglicherweise übrigen Mengen
erlauben, dass \(x\) sogar häufiger als in der Hälfte vorkommt.

Auch die Umkehrung gilt für endliche Mengen: Sind links \(r\)
und rechts \(s\) Mitglieder mit \(r\leq s\), nummeriert man
beide Klassen durch und ordnet das \(i\)-te Mitglied links
dem \(i\)-ten rechts zu. Diese Abbildung ist injektiv.
Für \(r=0\) übernimmt die leere Abbildung diese Aufgabe.
Auf dieses einfache Vergleichsprinzip greifen wir später
bei zwei Teilfamilien zurück.

Die allgemeine Frage lässt sich nun präzise formulieren:
\begin{quote}
  Gibt es im Träger jeder endlichen vereinigungsabgeschlossenen
  Familie \(M\) mit \(\bigcup M\neq\varnothing\) ein
  halbhäufiges Element?
\end{quote}
Die Bedingung an den Träger schließt die leere Familie und
die Familie \(\{\varnothing\}\) aus. In ihnen steht kein
Element zur Verfügung, dessen Häufigkeit die Behauptung
erfüllen könnte. Den Vereinigungsabschluss schreiben wir als
\[
  X,Y\in M\quad\Longrightarrow\quad X\cup Y\in M.
\]
Im Folgenden setzen wir diese Eigenschaft voraus und zeigen,
wie ein kleines Mitglied die gesuchte Zuordnung ermöglicht.

\par\Needspace{12\baselineskip}
\section*{Die Einermenge: Hinzufügen ohne Verlust}
\hypertarget{frankl.reading.singleton}{}

Nehmen wir zunächst an, dass \(\{a\}\in M\) ist. Für jede
Menge \(X\in M_a^-\) definieren wir
\[
  f(X)=X\cup\{a\}.
\]
Das Bild liegt in \(M\), weil \(X\) und \(\{a\}\)
Mitglieder von \(M\) sind und \(M\) unter Vereinigung
abgeschlossen ist. Außerdem enthält das Bild \(a\).
Wir erhalten daher eine wohldefinierte Abbildung
\[
  f\colon M_a^-\longrightarrow M_a^+.
\]
Die Abgeschlossenheit sichert, dass die Zuordnung innerhalb
der Familie bleibt. Für den Größenvergleich müssen wir
zusätzlich zeigen, dass sie keine zwei Mengen zusammenfallen
lässt.

Hier hilft gerade die Beschränkung auf Mengen ohne \(a\).
Für jedes \(X\in M_a^-\) gilt
\[
  (X\cup\{a\})\setminus\{a\}=X.
\]
Durch Entfernen des hinzugefügten Elements gewinnen wir die
ursprüngliche Menge vollständig zurück. Sind also
\(f(X)=f(Y)\) für \(X,Y\in M_a^-\), entfernen wir auf
beiden Seiten \(a\) und erhalten \(X=Y\). Damit ist
\(f\) injektiv. Nach \eqref{frankl-reading-half} folgt
\[
  |M_a^-|\leq|M_a^+|,
  \qquad |M_a^+|\geq\frac{|M|}{2}.
\]
Da \(\{a\}\) selbst zur Familie gehört, liegt \(a\) auch
im Träger. Somit ist \(a\) halbhäufig.
Der \FranklProofText{FranklSingletonMember}{Tabellenbeweis zum
Einermengenfall} führt die formale Ableitung aus.

Die Rolle der Voraussetzung ist genau erkennbar. Wir wissen
nicht etwa für jede vereinigungsabgeschlossene Familie, dass
\(X\cup\{a\}\) wieder in ihr liegt. Dafür brauchen wir
hier die Mitgliedschaft von \(\{a\}\). Unser Leitbeispiel
enthält keine Einermenge; auf dieses Beispiel lässt sich der
gerade bewiesene Fall also noch nicht anwenden.

Die Grenze \glqq mindestens die Hälfte\grqq{} kann schon im
Einermengenfall erreicht werden. In der Familie
\(\{\varnothing,\{a\}\}\) kommt \(a\) in genau einer von
zwei Mengen vor. Die Zuordnung verbindet die leere Menge mit
\(\{a\}\); es bleibt kein weiteres Mitglied auf der Zielseite.

\medskip\noindent\textbf{Wenn alle Mengen ein Element gemeinsam haben.}
\hypertarget{frankl.reading.intersection}{}
Ein noch direkterer Fall liegt vor, wenn \(M\neq\varnothing\)
und \(\bigcap M\neq\varnothing\) gelten. Wählen wir
\(a\in\bigcap M\), so enthält jedes Mitglied von \(M\)
dieses Element. Folglich ist \(M_a^-=\varnothing\) und
\(M_a^+=M\). Das Element \(a\) kommt sogar in allen Mengen
vor. Die leere Abbildung von \(M_a^-\) nach \(M_a^+\)
liefert auch hier den Injektionsbeweis. Für diesen Fall ist
der Vereinigungsabschluss gar nicht nötig. Die Beweistabellen
behandeln die \FranklProofText{FranklAvoidingEmpty}{leere
vermeidende Teilfamilie} und den
\FranklProofText{FranklNonemptyIntersection}{nichtleeren Durchschnitt}.

\par\Needspace{13\baselineskip}
\section*{Eine Zweiermenge führt zu vier Klassen}
\hypertarget{frankl.reading.partition}{}
\hypertarget{frankl.reading.classes}{}

Jetzt sei \(\{a,b\}\in M\) mit \(a\neq b\). Wir wollen
beweisen, dass mindestens eines der beiden Elemente
halbhäufig ist. Das Hinzufügen der Zweiermenge bleibt eine
zulässige Operation: Für \(X\in M\) liegt auch
\(X\cup\{a,b\}\) in \(M\).

Auf allen Mengen ohne \(a\) muss diese Abbildung jedoch
nicht injektiv sein. Die verschiedenen Mengen
\(\varnothing\) und \(\{b\}\) hätten beide das Bild
\(\{a,b\}\). Tatsächlich liegen diese Mengen beispielsweise
alle in der vereinigungsabgeschlossenen Potenzmenge von
\(\{a,b\}\). Dies ist ein anderes Beispiel; die Einermenge
\(\{b\}\) gehört nicht zu unserem Leitbeispiel.
Der Informationsverlust entsteht, weil nach dem Hinzufügen
nicht mehr erkennbar ist, ob \(b\) schon vorher vorhanden war.

Deshalb halten wir das Auftreten beider Elemente getrennt
fest. Wir bilden die vier Teilfamilien
\[
  \begin{aligned}
    M_{00}&=\{X\in M\mid a\notin X,\ b\notin X\},\\
    M_{01}&=\{X\in M\mid a\notin X,\ b\in X\},\\
    M_{10}&=\{X\in M\mid a\in X,\ b\notin X\},\\
    M_{11}&=\{X\in M\mid a\in X,\ b\in X\}.
  \end{aligned}
\]
Die erste Stelle des Index gehört zu \(a\), die zweite
zu \(b\). Eine Null bedeutet Abwesenheit, eine Eins
Anwesenheit. Jedes \(X\in M\) erhält genau eines dieser
vier Muster. Die Klassen sind somit paarweise disjunkt
und ergeben zusammen \(M\).

\FranklFourClassTable

Nun lassen sich beide Häufigkeitsfragen an derselben
Zerlegung ablesen:
\begin{equation}
  \begin{aligned}
    M_a^-&=M_{00}\mathbin{\dot\cup}M_{01},
    &M_a^+&=M_{10}\mathbin{\dot\cup}M_{11},\\
    M_b^-&=M_{00}\mathbin{\dot\cup}M_{10},
    &M_b^+&=M_{01}\mathbin{\dot\cup}M_{11}.
  \end{aligned}
  \tag{Z}\label{frankl-reading-partition}
\end{equation}
Der Punkt über dem Vereinigungszeichen erinnert an die
Disjunktheit. Zum Beispiel gehört eine Menge ohne \(a\)
entweder zu \(M_{00}\) oder zu \(M_{01}\), je nachdem,
ob sie \(b\) enthält. Beides zugleich ist unmöglich.
Die anderen drei Zerlegungen haben dieselbe Begründung.

Unsere Aufgabe hat jetzt zwei Teile. Zuerst vergleichen wir
die Klasse ohne beide Elemente mit der Klasse mit beiden.
Danach entscheiden die beiden gemischten Klassen, welches
Element wir als halbhäufig nachweisen können.

\par\Needspace{12\baselineskip}
\section*{Zwei Zuordnungen ergänzen sich zum Beweis}
\hypertarget{frankl.reading.injection}{}

Auf \(M_{00}\) ist das Hinzufügen beider Elemente sicher.
Definieren wir
\[
  \varphi\colon M_{00}\longrightarrow M_{11},
  \qquad \varphi(X)=X\cup\{a,b\},
\]
so liegt das Bild wegen des Vereinigungsabschlusses in \(M\)
und enthält \(a\) und \(b\). Die Abbildung hat also
tatsächlich den angegebenen Zielbereich. Für \(X\in M_{00}\)
enthält \(X\) keines der beiden Elemente, daher gilt
\[
  \varphi(X)\setminus\{a,b\}=X.
\]
Aus \(\varphi(X)=\varphi(Y)\) folgt durch Entfernen
beider Elemente \(X=Y\). Damit ist \(\varphi\) injektiv
und wir haben den grundlegenden Vergleich
\begin{equation}
  |M_{00}|\leq|M_{11}|.
  \tag{I}\label{frankl-reading-injection}
\end{equation}
Nun sind \(M_{01}\) und \(M_{10}\) endliche Teilfamilien
von \(M\). Ihre Anzahlen lassen sich vergleichen.
Wir behandeln beide Möglichkeiten und bauen in jedem Fall
die vollständige benötigte Injektion.

\medskip\noindent\textbf{Erster Fall: \(|M_{01}|\leq|M_{10}|\).}
Nach dem Vergleichsprinzip für endliche Mengen gibt es eine
Injektion \(g\colon M_{01}\to M_{10}\). Sie braucht keine
Vereinigungsabbildung zu sein: Wir können die kleinere Klasse
durchnummerieren und ihren Mitgliedern verschiedene Mitglieder
der größeren Klasse zuweisen.
Zusammen mit \(\varphi\) definieren wir
\[
  F\colon M_a^-\longrightarrow M_a^+,
  \qquad
  F(X)=
  \begin{cases}
    \varphi(X),&X\in M_{00},\\
    g(X),&X\in M_{01}.
  \end{cases}
\]
Die beiden Fälle sind disjunkt und decken nach
\eqref{frankl-reading-partition} den ganzen Definitionsbereich
ab. Jeder Wert ist eindeutig bestimmt und liegt in \(M_a^+\).

\FranklInjectionDiagram

Für die Injektivität sind drei Situationen zu prüfen.
Liegen zwei Ausgangsmengen in \(M_{00}\), folgt ihre
Gleichheit aus der Injektivität von \(\varphi\).
Liegen beide in \(M_{01}\), folgt sie aus der Injektivität
von \(g\). Liegt eine Menge in jeder der beiden Klassen,
so können ihre Bilder gar nicht gleich sein: Das Bild des
ersten Zweigs gehört zu \(M_{11}\) und enthält \(b\),
das Bild des zweiten gehört zu \(M_{10}\) und enthält
\(b\) nicht. Somit ist \(F\) injektiv.

Es folgt \(|M_a^-|\leq|M_a^+|\), und nach
\eqref{frankl-reading-half} ist \(a\) halbhäufig.
Dies ist der \FranklProofText{FranklPairPositiveBranch}{positive
Zweig des Tabellenbeweises}.

\medskip\noindent\textbf{Zweiter Fall: \(|M_{01}|>|M_{10}|\).}
Jetzt gibt es eine Injektion
\(h\colon M_{10}\to M_{01}\). Wir wechseln zum Element
\(b\) und setzen
\[
  G\colon M_b^-\longrightarrow M_b^+,
  \qquad
  G(X)=
  \begin{cases}
    \varphi(X),&X\in M_{00},\\
    h(X),&X\in M_{10}.
  \end{cases}
\]
Wieder decken die disjunkten Ausgangsklassen nach
\eqref{frankl-reading-partition} genau \(M_b^-\) ab.
Die Bilder des ersten Zweigs liegen in \(M_{11}\), die
des zweiten in \(M_{01}\), also sämtlich in \(M_b^+\).
Innerhalb eines Zweigs ist die Zuordnung injektiv, weil
\(\varphi\) beziehungsweise \(h\) injektiv ist.
Bilder aus verschiedenen Zweigen sind ebenfalls verschieden:
Oben enthalten sie \(a\), unten enthalten sie \(a\) nicht.
Damit ist \(G\) injektiv und \(b\) halbhäufig.
Der \FranklProofText{FranklPairNegativeBranch}{zweite Tabellenzweig}
führt diesen Nachweis aus.

Die beiden Fälle erfassen alle Möglichkeiten. Wir haben
deshalb vollständig bewiesen:
\hypertarget{frankl.reading.pair}{}
\begin{quote}
  Enthält eine endliche vereinigungsabgeschlossene Familie
  die Zweiermenge \(\{a,b\}\), so kommt mindestens eines
  der Elemente \(a,b\) in mindestens der Hälfte ihrer
  Mitglieder vor.
\end{quote}
Beide Elemente gehören zum Träger, weil \(\{a,b\}\in M\)
ist. Den Abschluss bildet der
\FranklProofText{FranklFamPairMemberFrankl}{Tabellenbeweis zum
Zweiermengenfall}; seine Aussage bleibt in Band~46.

Im Leitbeispiel sind die gemischten Klassen gleich groß.
Wir können \(g(\{b,c\})=\{a,c\}\) wählen und erhalten
\[
  F(\varnothing)=\{a,b\},\qquad
  F(\{b,c\})=\{a,c\}.
\]
Dies sind zwei verschiedene Mengen mit \(a\). Die dritte,
\(\{a,b,c\}\), wird gar nicht benötigt. So erklärt die
Zuordnung die anfangs beobachteten drei Vorkommen von \(a\)
bei nur zwei Mengen ohne \(a\). Wegen der Gleichheit der
Mischklassen kann man hier auch in der anderen Richtung
zuordnen und \(b\) ebenso nachweisen.

\par\Needspace{12\baselineskip}
\section*{Derselbe Beweis als Zählgleichung}
\hypertarget{frankl.reading.counting}{}

Die beiden zusammengesetzten Injektionen zeigen genau, woher
der Häufigkeitsvergleich kommt. Man kann ihren gemeinsamen
Kern auch in einer kurzen Rechnung zusammenfassen.
Schreiben wir \(n_{ij}=|M_{ij}|\), so ergeben sich
\[
  |M|=n_{00}+n_{01}+n_{10}+n_{11},\qquad
  |M_a^+|=n_{10}+n_{11},\quad
  |M_b^+|=n_{01}+n_{11}.
\]
Die Summe der beiden Vorkommenszahlen ist daher
\begin{equation}
  \begin{aligned}
    |M_a^+|+|M_b^+|
      &=n_{10}+n_{01}+2n_{11}\\
      &=|M|+(n_{11}-n_{00})\\
      &\geq|M|.
  \end{aligned}
  \tag{S}\label{frankl-reading-sum}
\end{equation}
Die letzte Zeile ist genau der Vergleich
\eqref{frankl-reading-injection}. Wären beide Vorkommenszahlen
kleiner als \(|M|/2\), wäre ihre Summe kleiner als \(|M|\).
Das widerspricht \eqref{frankl-reading-sum}. Mindestens eine
von ihnen erreicht also die geforderte Hälfte.

In dieser Rechnung zählt jedes Mitglied von \(M_{11}\)
zweimal, weil es beide Elemente enthält. Jedes Mitglied
einer gemischten Klasse zählt einmal, eines von \(M_{00}\)
gar nicht. Gegenüber dem einmaligen Zählen aller Mitglieder
gleicht die zusätzliche Zählung in \(M_{11}\) somit die
fehlende Zählung in \(M_{00}\) aus. Genau diesen Ausgleich
sichert die Injektion \(\varphi\).

Im Leitbeispiel sind
\[
  (n_{00},n_{01},n_{10},n_{11})=(1,1,1,2),
  \qquad |M_a^+|+|M_b^+|=3+3=5+(2-1).
\]
Der Überschuss in der Klasse mit beiden Elementen wird so
auch zahlenmäßig sichtbar. Welche einzelne Vorkommenszahl
größer ist, entscheidet dagegen der Vergleich der Mischklassen:
\[
  |M_a^+|-|M_b^+|=n_{10}-n_{01}.
\]
Damit passt die kurze Rechnung genau zur vorherigen Fallwahl.
Sind die gemischten Klassen gleich groß, sind beide Elemente
halbhäufig. Sind sie verschieden groß, beweist unser Verfahren
die Behauptung für das häufigere der beiden Elemente.

\par\Needspace{13\baselineskip}
\section*{Wie weit das Ergebnis reicht}
\hypertarget{frankl.reading.scope}{}

Der Zweiermengenfall besagt mehr als nur, dass es irgendwo
im Träger ein halbhäufiges Element gibt: Eines der beiden
Elemente der enthaltenen Zweiermenge genügt bereits.
Andere Mitglieder der Familie dürfen beliebig viele weitere
Elemente besitzen. Der Beweis verlangt auch nicht, dass die
leere Menge zur Familie gehört. Falls \(M_{00}\) leer ist,
bleibt \(\varphi\) als leere Injektion gültig; alle
anschließenden Argumente funktionieren unverändert.

Ebenso wenig wurde die Endlichkeit der einzelnen Mitglieder
benötigt. Das Hinzufügen und anschließende Entfernen von
\(a,b\) funktioniert auch für unendliche Mengen.
Endlich sein muss in dieser Darstellung die \emph{Familie}:
Wir zählen ihre Mitglieder und vergleichen die beiden
gemischten Klassen durch ihre endlichen Anzahlen.
Die Schreibweise \(|M|/2\) verwendet diesen endlichen
Zählbegriff.

Was geschieht, wenn eine Dreiermenge \(A=\{a,b,c\}\)
als Mitglied vorliegt? Wir fassen die Mengen jetzt danach
zusammen, wie viele Elemente von \(A\) sie enthalten:
\[
  \begin{array}{c|cccc}
    \text{Anzahl enthaltener Elemente von }A & 0 & 1 & 2 & 3\\
    \text{Anzahl solcher Mengen in }M       & n_0 & n_1 & n_2 & n_3
  \end{array}
\]
Dabei zählt etwa \(n_1\) alle Mengen, die genau eines von
\(a,b,c\) enthalten, gleichgültig welches. Jedes Mitglied
gehört genau zu einer der vier Gruppen, also gilt
\(|M|=n_0+n_1+n_2+n_3\).

Addieren wir die Vorkommenszahlen von \(a,b,c\), so trägt
eine Menge ohne diese Elemente nichts bei. Eine Menge mit
genau einem trägt einmal bei, eine mit genau zwei zweimal,
eine mit allen dreien dreimal. Deshalb gilt
\[
  |M_a^+|+|M_b^+|+|M_c^+|=0n_0+1n_1+2n_2+3n_3.
\]
Um über diese Summe eine Vorkommenszahl von mindestens
\(|M|/2\) zu erzwingen, genügte eine Summe von mindestens
\(3|M|/2\). Wir setzen die Zerlegung von \(|M|\) ein
und fassen die Koeffizienten zusammen:
\[
  \begin{aligned}
    n_1+2n_2+3n_3-\frac32|M|
      &=n_1+2n_2+3n_3-\frac32(n_0+n_1+n_2+n_3)\\
      &=-\frac32n_0-\frac12n_1+\frac12n_2+\frac32n_3\\
      &=\frac12(n_2-n_1)+\frac32(n_3-n_0).
  \end{aligned}
\]
Erneut ist das Hinzufügen von \(A\) auf den Mengen ohne
\(a,b,c\) injektiv: Entfernen wir \(A\) wieder, erhalten
wir die Ausgangsmenge. Die Bilder enthalten alle drei
Elemente. Das liefert \(n_3\geq n_0\), sagt aber noch
nichts über \(n_2-n_1\). Der zweite Vergleich fehlt.
Das nächste Beispiel zeigt, dass sogar alle drei Elemente
von \(A\) seltener als in der Hälfte vorkommen können.

\par\Needspace{16\baselineskip}
\section*{Eine Dreiermenge ohne halbhäufiges Element}
\hypertarget{frankl.reading.rare-triple}{}
\hypertarget{frankl.reading.triple-example}{}

Es gibt eine vereinigungsabgeschlossene Familie mit einem
Dreiermitglied \(A=\{a,b,c\}\), in der jedes dieser drei
Elemente nur in neun von neunzehn Mengen vorkommt. Die Familie
hat insgesamt sieben Trägerelemente: Außer \(a,b,c\) treten
\(d,e,f,g\) auf. Diese vier weiteren Elemente sind häufig.
Damit scheitert die direkte Übertragung des Zweiermengensatzes
auf ein vorgegebenes Dreiermitglied; Frankls allgemeine
Forderung nach irgendeinem halbhäufigen Element bleibt in
diesem Beispiel erfüllt. Die folgende Konstruktion stammt
von van der Hout und Roos; wir gruppieren ihre Inzidenzmatrix
so, dass der Vereinigungsabschluss unmittelbar prüfbar wird.\footnote{Van der
Hout und Roos, \href{https://doi.org/10.23952/jano.8.2026.1.02}{\emph{Some results
and conjectures related to Frankl's union closed conjecture}},
Abschnitt~4, Abbildung~1, S.~13--14. Dort heißen die sieben
Elemente \(1,\ldots,7\); hier verwenden wir \(a,\ldots,g\).}

Alle sieben Elemente seien verschieden. Wir setzen
\[
  \begin{aligned}
    A&=\{a,b,c\}, & B&=\{d,e,f,g\},\\
    D_a&=\{d,e,g\}, &D_b&=\{d,f,g\}, &D_c&=\{e,f,g\}.
  \end{aligned}
\]
Unsere Familie besteht aus diesen Bausteinen:
\[
  \begin{aligned}
    M={}&\{\varnothing,A\}\\
       &{}\cup\bigcup_{x\in A}\{D_x,\ D_x\cup\{x\},\ D_x\cup A\}\\
       &{}\cup\{B\cup T\mid T\subseteq A\}.
  \end{aligned}
\]
Der erste Baustein enthält zwei Mengen. Es folgen drei
Blöcke mit je drei Mengen. Im letzten Baustein ergänzen wir
\(B\) um jede der acht Teilmengen von \(A\). Die Bausteine
sind disjunkt: Ihr Anteil in \(B\) ist entweder leer,
eines der drei verschiedenen \(D_x\) oder ganz \(B\).
Auch innerhalb jedes Bausteins sind die Mengen verschieden.
Somit gilt \(|M|=2+3\cdot3+8=19\). Die kleinsten
nichtleeren Mitglieder sind Dreiermengen; Einermengen und
Zweiermengen kommen nicht vor.

\FranklRareTripleTable

\medskip\noindent\textbf{Warum die Familie vereinigungsabgeschlossen ist.}
Die leere Menge ändert bei einer Vereinigung nichts.
Innerhalb eines Dreierblocks sind die Mengen ineinander
enthalten:
\[
  D_x\subseteq D_x\cup\{x\}\subseteq D_x\cup A.
\]
Ihre Vereinigung ist daher die größere der beiden Mengen.
Für verschiedene \(x,y\in A\) gilt dagegen
\(D_x\cup D_y=B\). Die Vereinigung zweier Mengen aus
verschiedenen Blöcken hat folglich die Form \(B\cup T\)
mit \(T\subseteq A\) und liegt im letzten Baustein.
Vereinigen wir \(A\) mit einer Menge eines Dreierblocks,
entsteht \(D_x\cup A\), das oberste Mitglied dieses Blocks.
Schließlich bleibt jede Vereinigung mit \(B\cup T\)
im letzten Baustein: Sie enthält weiterhin ganz \(B\)
und ergänzt nur weitere Elemente von \(A\).
Zusammen mit \(A\cup A=A\) erfasst dies alle Fälle.

\par\Needspace{10\baselineskip}
\medskip\noindent\textbf{Warum alle drei Elemente von \(A\) selten sind.}
Ein festes \(x\in A\) kommt in \(A\) selbst vor,
in genau einem mittleren Blockmitglied \(D_x\cup\{x\}\)
und in allen drei oberen Blockmitgliedern \(D_y\cup A\).
Unter den acht Mengen \(B\cup T\) enthalten genau vier
das Element \(x\): Die beiden anderen Elemente von \(A\)
können unabhängig hinzukommen oder fehlen. Also gilt
\[
  |M_a^+|=|M_b^+|=|M_c^+|=1+1+3+4=9<\frac{19}{2}.
\]
Die rechte Hälfte der Tafel zeigt zugleich
\(|M_d^+|=|M_e^+|=|M_f^+|=14\) und \(|M_g^+|=17\).
Halbhäufige Elemente sind somit vorhanden, liegen aber
sämtlich außerhalb des ausgewählten Dreiermitglieds.

Auch die zuvor erklärte Zählgleichung wird konkret:
\[
  (n_0,n_1,n_2,n_3)=(5,6,3,5).
\]
Ohne ein Element von \(A\) sind \(\varnothing\), die
drei \(D_x\) und \(B\). Genau eines enthalten die
drei mittleren Blockmitglieder und die drei Mengen
\(B\cup\{x\}\). Genau zwei enthalten die drei Mengen
\(B\cup\{a,b\}\), \(B\cup\{a,c\}\),
\(B\cup\{b,c\}\). Alle drei enthalten \(A\), die
drei oberen Blockmitglieder und \(B\cup A\). Daher
\[
  \begin{aligned}
    9+9+9&=6+2\cdot3+3\cdot5=27<\frac32\cdot19,\\
    \frac12(3-6)+\frac32(5-5)&=-\frac32.
  \end{aligned}
\]
Hier sind \(n_3\) und \(n_0\) sogar gleich groß.
Der Fehlbetrag entsteht vollständig dadurch, dass mehr
Mengen genau eines als genau zwei der drei Elemente enthalten.

Die \FranklMainLink{} führt zu den Aussagen in Band~46.
Die \FranklProofsLink{} enthalten die 17 bisherigen
Ableitungen zu den Einermengen, zum gemeinsamen Durchschnitt
und zu den Zweiermengen sowie die
\FranklProofText{FranklRareTripleExists}{formale Herleitung des
neuen Beispiels}. Die Lesefassung erklärt die Beweise selbst;
die Tabellen halten ihre logischen Schritte fest.

\par\Needspace{16\baselineskip}
\section*{Historischer Abriss}
\hypertarget{frankl.reading.history}{}

Die Vermutung geht auf Péter Frankl zurück und wird auf
1979 datiert. Rein van der Hout und Kees Roos verweisen
hierfür auf Frankls Beitrag \emph{Extremal set systems}
im \emph{Handbook of Combinatorics} von 1995. Dort begegnet
die Frage in der dualen Form für durchschnittsabgeschlossene
Familien.\footnote{Rein van der Hout und Kees Roos,
\href{https://doi.org/10.23952/jano.8.2026.1.02}{\emph{Some results and conjectures related to Frankl's union closed conjecture}},
\emph{Journal of Applied and Numerical Optimization}~8 (2026),
S.~11--18, Einleitung und Bemerkung~1.1. Die Datierung und
der Verweis auf Frankls Beitrag werden hier nach dieser
Quelle wiedergegeben.}

Die Verbindung beider Formen lässt sich direkt verstehen.
Ersetzt man jedes Mitglied \(X\) einer Familie auf einem
Träger \(U\) durch sein Komplement \(U\setminus X\),
so wird die Vereinigung durch die De-Morgansche Regel zum
Durchschnitt:
\[
  U\setminus(X\cup Y)
    =(U\setminus X)\cap(U\setminus Y).
\]
Dabei kommt ein festes Element genau in denjenigen
Komplementmengen vor, deren ursprüngliche Mengen es nicht
enthalten. \glqq Mindestens die Hälfte\grqq{} auf der einen
Seite entspricht daher \glqq höchstens die Hälfte\grqq{}
auf der anderen. Hinter beiden Formulierungen steht dieselbe
Häufigkeitsfrage.

Die hier bewiesenen kleinen Fälle gehören zur frühen
Entwicklung des Problems. Van der Hout und Roos verweisen
für den Zweiermengenfall auf einen Satz von D.~G.~Sarvate
und J.-C.~Renaud aus dem Jahr 1989. In ihrem eigenen
Abschnitt~3 führen sie die beiden Spezialfälle erneut vor;
die Injektion von den Mengen ohne beide Elemente zu den
Mengen mit beiden bildet auch dort den entscheidenden
Schritt.\footnote{Van der Hout und Roos, a.~a.~O.,
Abschnitt~3 und Beginn von Abschnitt~4. Dort wird auf
D.~G.~Sarvate und J.-C.~Renaud, \emph{On the union-closed
sets conjecture}, \emph{Ars Combinatoria}~27 (1989),
S.~149--153, Satz~2 verwiesen.}

Einen anderen Zugang eröffnete Justin Gilmer 2022 mit
Methoden der Informationstheorie. Er bewies für die
allgemeine endliche Fragestellung die Existenz eines
Elements, das in mindestens einem Prozent der Mengen
vorkommt. Die Konstante hängt weder von der Größe des
Trägers noch von der Zahl der Familienmitglieder ab.
Sein Argument untersucht zufällig gewählte Mengen und
die Information in ihrer Vereinigung.\footnote{Justin Gilmer,
\href{https://arxiv.org/abs/2211.09055v2}{\emph{A constant lower bound for the union-closed sets conjecture}},
arXiv:2211.09055, 2022, Fassung vom 28.~November 2022.
Die genannte Ein-Prozent-Schranke ist das Ergebnis dieser
Arbeit und keine Angabe der jeweils besten späteren Schranke.}
Die volle Hälfte der Spezialfälle und die kleinere
allgemeine Schranke beantworten unterschiedliche Fragen:
Hier nutzen wir ein bestimmtes kleines Mitglied; eine
allgemeine Methode muss auch ohne ein solches Mitglied
auskommen. An beiden Zugängen wird sichtbar, wie eng
Vereinigung und Häufigkeit miteinander verbunden sind.

\par\Needspace{16\baselineskip}
\section*{Schlussbemerkung}
\hypertarget{frankl.reading.conclusion}{}

In unserer Familie aus fünf Mengen kamen \(a,b,c\) jeweils
dreimal vor. Die Tafel machte das Ergebnis sichtbar;
der Beweis erklärt, weshalb eine enthaltene Einermenge
oder Zweiermenge mindestens die Hälfte garantiert.
Vereinigungsabschluss lässt uns ein kleines Mitglied
hinzufügen, ohne die Familie zu verlassen. Beschränken
wir den Ausgangsbereich passend, lässt sich die ursprüngliche
Menge anschließend zurückgewinnen. Aus diesem fehlenden
Informationsverlust entsteht die Injektivität, aus der
Injektivität der Vergleich der Anzahlen.

Bei einer Einermenge genügt eine einzige Zuordnung. Bei
einer Zweiermenge trennt die Vierfeldertafel die sichere
Zuordnung von den beiden gemischten Klassen. Deren Vergleich
entscheidet, welches der zwei Elemente die Hälfte erreicht.
Die disjunkten Zielklassen sichern, dass die zusammengesetzte
Abbildung injektiv bleibt. Die kurze Zählgleichung fasst
denselben Ausgleich in Zahlen zusammen.

Das Beispiel mit neunzehn Mengen zeigt die genaue Grenze
dieser Aussage: Alle drei Elemente eines enthaltenen
Dreiermitglieds können unterhalb der Hälfte bleiben.
Die übrigen Trägerelemente erfüllen dort Frankls Forderung.
So erklärt die Lesefassung zugleich den Erfolg bei ein
oder zwei Elementen und das Scheitern derselben lokalen
Behauptung bei drei Elementen. Eine geeignete Zerlegung
verwandelt die Regel über Vereinigungen in einen Vergleich
der Häufigkeiten; das Gegenbeispiel zeigt, welchen Vergleich
der Vereinigungsabschluss allein nicht erzwingt.
